\documentclass[12pt, a4paper]{article}

\usepackage[utf8]{inputenc}
\usepackage[T1]{fontenc}
\usepackage{geometry}
\usepackage{setspace}
\usepackage{titlesec}
\usepackage{hyperref}
\usepackage{graphicx}
\usepackage{longtable}
\usepackage{array}
\usepackage{booktabs}
\usepackage{tabularx}

\geometry{margin=1in}
\doublespacing

\hypersetup{
    colorlinks=true,
    linkcolor=blue,
    citecolor=blue,
    urlcolor=blue
}

\title{%
    \vspace{2cm}
    \textbf{Evaluation Protocol (Sprint 4 Bridge)} \\
    \vspace{0.5cm}
    \large CISC-727: Advanced Research Explorations (ARE) - I \\
    \vspace{0.3cm}
    \large Assignment A03
}
\author{%
    Kenneth Peter Fernandes \\
    \vspace{0.3cm}
    Harrisburg University of Science and Technology
}
\date{Spring 2026}

\begin{document}

% ============================================================
% Title Page
% ============================================================
\begin{titlepage}
    \centering
    \vspace*{2cm}
    {\LARGE \textbf{Evaluation Protocol (Sprint 4 Bridge)} \par}
    \vspace{1cm}
    {\large CISC-727: Advanced Research Explorations (ARE) - I \par}
    \vspace{0.5cm}
    {\large Assignment A03 \par}
    \vspace{1.5cm}
    {\large \textit{FlashAttention Speedup Under Multi-GPU Communication Overhead: DDP vs FSDP} \par}
    \vspace{1.5cm}
    {\Large Kenneth Peter Fernandes \par}
    \vspace{0.5cm}
    {\large Harrisburg University of Science and Technology \par}
    \vspace{1cm}
    {\large Spring 2026 \par}
\end{titlepage}

% ============================================================
% Section A: Evaluation Protocol
% ============================================================
\section{Evaluation Protocol}

% ============================================================
\subsection{Research Question + Hypothesis}

\textbf{A02 Gap (Continuity):} The literature on efficient attention demonstrates that FlashAttention achieves significant speedups over standard attention through IO-aware kernel optimization on single-GPU settings. However, it remains underexplored how this compute-side advantage interacts with multi-GPU communication overhead---specifically whether FlashAttention's relative speedup is preserved, amplified, or diminished as the communication burden increases across different parallelization strategies.

\vspace{1em}
\textbf{Research Question:} Does FlashAttention's speedup over standard attention diminish as communication overhead increases when scaling from single-GPU to multi-GPU training using DDP and FSDP?

\vspace{1em}
\textbf{Hypothesis:} FlashAttention's relative speedup over standard attention diminishes as communication overhead increases. Specifically:
\begin{itemize}
    \item On a single GPU (no communication), FlashAttention will show its maximum relative speedup over standard attention.
    \item Under DDP (gradient all-reduce at the end of backward pass---low communication), the speedup will be partially preserved because communication overhead is modest relative to compute time.
    \item Under FSDP (all-gather per layer + gradient reduce-scatter---high communication), the speedup will shrink further because communication becomes a larger fraction of total step time, reducing the relative impact of compute-side IO-aware optimization.
\end{itemize}

\vspace{1em}
\textbf{Falsification Criteria:} The hypothesis is refuted if FlashAttention's relative speedup (\%) remains constant or increases when moving from single-GPU $\rightarrow$ DDP $\rightarrow$ FSDP. Specifically, if the speedup under FSDP is equal to or greater than the speedup on single GPU (within measurement noise), the hypothesis is rejected.

% ============================================================
\subsection{System Under Test}

\begin{singlespace}
\small
\begin{tabularx}{\textwidth}{>{\raggedright\arraybackslash}p{4cm} X}
\toprule
\textbf{Component} & \textbf{Specification} \\
\midrule
Platform & Kaggle Notebooks \\
GPUs & 2$\times$ NVIDIA Tesla T4 (16 GB HBM2 each) \\
Interconnect & PCIe (not NVLink) \\
Framework & PyTorch 2.x \\
Distributed Backend & NCCL \\
Model & ViT-Small (patch size 16, embed dim 384, 6 heads, 12 layers) \\
Dataset & CIFAR-10 (50,000 train images, 10 classes, 32$\times$32 pixels) \\
Batch Size & 64 per GPU (held constant across all runs) \\
Training Steps & 100 steps (pilot; sufficient for timing, not convergence) \\
Attention Backends & Standard attention (math backend) vs FlashAttention (via \texttt{torch.nn.functional.scaled\_dot\_product\_attention} with flash backend) \\
Precision & float32 (default); float16 via AMP (when toggled) \\
Random Seed & 42 (fixed across all runs) \\
\bottomrule
\end{tabularx}
\end{singlespace}

% ============================================================
\subsection{Baselines}

\textbf{Baseline \#1---Single GPU, Standard Attention (R1):}
Single T4 GPU training with PyTorch's default math-based attention. This is the absolute reference point---no parallelism, no IO-aware optimization. All speedups are measured relative to this.

\vspace{1em}
\textbf{Baseline \#2---Single GPU, FlashAttention (R2):}
Single T4 GPU training with FlashAttention enabled via \texttt{scaled\_dot\_product\_attention}. This establishes the maximum FlashAttention speedup achievable without any communication overhead. The difference between R1 and R2 isolates FlashAttention's compute-side benefit.

% ============================================================
\subsection{Independent Variables (Factors)}

\begin{enumerate}
    \item \textbf{Attention Backend} (2 levels):
    \begin{itemize}
        \item Standard attention (PyTorch math backend)
        \item FlashAttention (PyTorch flash backend via SDPA)
    \end{itemize}

    \item \textbf{Parallelization Strategy} (3 levels):
    \begin{itemize}
        \item Single GPU (no parallelism)
        \item DDP (DistributedDataParallel---low communication: all-reduce after backward)
        \item FSDP (FullyShardedDataParallel---high communication: all-gather per layer + reduce-scatter)
    \end{itemize}

    \item \textbf{Compute-side Optimization} (2 levels):
    \begin{itemize}
        \item Off (float32)
        \item Mixed Precision / AMP (float16 via \texttt{torch.autocast})
    \end{itemize}

    \item \textbf{Comm-side Optimization} (2 levels):
    \begin{itemize}
        \item Off (default sequential compute then communicate)
        \item Overlap Compute/Comm (overlap gradient communication with backward pass computation)
    \end{itemize}
\end{enumerate}

\vspace{1em}
\textbf{Why DDP + FSDP:}
These two strategies share the same data-parallel philosophy (each GPU processes different data) but differ fundamentally in communication pattern. DDP communicates once at the end (all-reduce gradients), while FSDP communicates every layer (all-gather parameters + reduce-scatter gradients). This creates a natural communication overhead gradient (low $\rightarrow$ high) that directly tests whether FlashAttention's compute-side benefit survives increasing communication load. Model parallelism and pipeline parallelism were rejected because they introduce orthogonal bottlenecks (tensor splitting, pipeline bubbles) that would confound the attention-kernel-focused hypothesis.

\vspace{1em}
\textbf{Why AMP:}
Both AMP and FlashAttention optimize memory access---AMP at the data-type level (float16 reduces memory bandwidth), FlashAttention at the kernel level (tiling reduces HBM traffic). Testing AMP reveals whether these two memory-access optimizations stack independently or show diminishing returns when combined.

\vspace{1em}
\textbf{Why Overlap Compute/Comm:}
FlashAttention reduces compute time per step. Overlap compute/comm hides communication behind compute. If FlashAttention shortens compute, there is less time available to hide communication---making overlap potentially less effective with FlashAttention than with standard attention. This is a testable interaction.

% ============================================================
\subsection{Dependent Variables (Metrics)}

\textbf{Required Metrics:}
\begin{itemize}
    \item \textbf{Time per step} (seconds): wall-clock time for one training step (forward + backward + optimizer), measured via \texttt{torch.cuda.Event} for precision.
    \item \textbf{Throughput} (samples/second): \texttt{batch\_size / time\_per\_step}.
    \item \textbf{Speedup}: \texttt{time\_per\_step(baseline) / time\_per\_step(variant)}, relative to R1.
    \item \textbf{Scaling efficiency}: \texttt{speedup / number\_of\_GPUs} (ideal = 1.0).
\end{itemize}

\vspace{1em}
\textbf{Additional Metrics:}
\begin{itemize}
    \item \textbf{Compute vs communication breakdown}: separate timing of forward pass, backward pass, and communication (gradient sync) using CUDA events. This directly measures whether communication is the bottleneck that erodes FlashAttention's advantage.
    \item \textbf{FlashAttention relative speedup (\%)}: \texttt{(time\_standard $-$ time\_flash) / time\_standard $\times$ 100}, computed within each parallelization strategy. This is the primary metric for testing the hypothesis---if this percentage drops from single-GPU $\rightarrow$ DDP $\rightarrow$ FSDP, the hypothesis is supported.
\end{itemize}

% ============================================================
\subsection{Ablation Plan}

Ablations follow a one-factor-at-a-time discipline. Each run changes exactly one variable from a reference configuration, so the effect of each factor can be isolated.

\vspace{1em}
\textbf{Stage 1---Attention Backend $\times$ Parallelization (core hypothesis test):}

Hold constant: batch size, model, dataset, steps, precision (float32), no optimizations.
\begin{itemize}
    \item R1: Single GPU + Standard Attention (absolute baseline)
    \item R2: Single GPU + FlashAttention
    \item R3: DDP + Standard Attention
    \item R4: DDP + FlashAttention
    \item R5: FSDP + Standard Attention
    \item R6: FSDP + FlashAttention
\end{itemize}

Compare FlashAttention speedup: R2 vs R1, R4 vs R3, R6 vs R5. If speedup \% decreases across these pairs, hypothesis is supported.

\vspace{1em}
\textbf{Stage 2---Optimization Ablations (on DDP + FlashAttention as reference):}

Hold constant: DDP, FlashAttention, batch size, model, dataset, steps.
\begin{itemize}
    \item R4: DDP + FlashAttention (reference, no optimizations)
    \item R7: DDP + FlashAttention + AMP only
    \item R8: DDP + FlashAttention + Overlap Compute/Comm only
    \item R9: DDP + FlashAttention + AMP + Overlap (combined)
\end{itemize}

This tests whether each optimization provides independent benefit and whether they stack.

% ============================================================
\subsection{Acceptance Criteria}

\begin{enumerate}
    \item \textbf{Warmup policy:} The first 10 steps of each run are discarded. Timing is measured over steps 11--100 (90 steps). This avoids CUDA warmup, JIT compilation, and data loader initialization artifacts.
    \item \textbf{Repetitions:} Each run is repeated 3 times. The \textbf{median} of the 3 runs is reported to reduce outlier sensitivity.
    \item \textbf{Convergence constraint:} Since this is a pilot timing study (100 steps), convergence is not required. However, training loss must decrease over the 100 steps to confirm the model is learning and the setup is not broken.
    \item \textbf{Success threshold:} A speedup or slowdown is considered meaningful only if it exceeds \textbf{5\%} relative to the reference. Differences within 5\% are reported as ``within noise.''
    \item \textbf{Hardware consistency:} All runs must execute on the same Kaggle session (same GPU allocation) to avoid hardware variability across sessions.
\end{enumerate}

% ============================================================
\subsection{Threats to Validity}

\begin{singlespace}
\small
\begin{longtable}{>{\raggedright\arraybackslash}p{3.5cm} >{\raggedright\arraybackslash}p{4.5cm} >{\raggedright\arraybackslash}p{5.5cm}}
\toprule
\textbf{Threat} & \textbf{Impact} & \textbf{Mitigation} \\
\midrule
\endfirsthead
\toprule
\textbf{Threat} & \textbf{Impact} & \textbf{Mitigation} \\
\midrule
\endhead
\bottomrule
\endfoot

Kaggle PCIe interconnect & T4s connect via PCIe (not NVLink), inflating communication overhead compared to production setups & Acknowledged as a limitation; results may overstate communication impact. Report interconnect type explicitly. \\
\midrule
Kaggle session variability & Different sessions may allocate different physical GPUs or share host resources & Run all experiments in a single session. Report if session restarts occur. \\
\midrule
Warmup / JIT effects & First few steps are slower due to CUDA kernel compilation and caching & Discard first 10 steps; report only steps 11--100. \\
\midrule
Small model on 2 GPUs & ViT-Small may not stress communication enough; compute may dominate even under FSDP & Acknowledged. If comm overhead is negligible, results will show flat speedup across strategies (which would refute the hypothesis---a valid outcome). \\
\midrule
Batch size changes & Different strategies may require different batch sizes for memory reasons & Batch size is held constant (64/GPU). If FSDP requires adjustment, it will be documented and justified. \\
\midrule
CIFAR-10 image size & 32$\times$32 images produce short sequences (4$\times$4 = 16 patches with patch size 16), which may not stress attention enough & Acknowledged. Sequence length may be too short for FlashAttention to show significant advantage. If so, patch size can be reduced to increase sequence length (e.g., patch size 4 $\rightarrow$ 64 patches). \\
\midrule
Random seed sensitivity & Single seed may produce unrepresentative training dynamics & Fixed seed (42) for reproducibility. Timing metrics (not accuracy) are the primary outcome, reducing seed sensitivity. \\
\midrule
Mixed precision numerical effects & AMP may change convergence behavior or introduce instability & Loss values are logged to verify training stability. AMP runs are compared on timing, not accuracy. \\
\midrule
Torch compile / backend caching & PyTorch may cache kernels across runs within the same process & Each run configuration is executed in a fresh process or after cache clearing. \\

\end{longtable}
\end{singlespace}

% ============================================================
% Section B: Run Plan Table
% ============================================================
\section{Run Plan Table}

\begin{singlespace}
\small
\begin{longtable}{>{\raggedright\arraybackslash}p{0.8cm} >{\raggedright\arraybackslash}p{1.8cm} >{\raggedright\arraybackslash}p{1.8cm} >{\raggedright\arraybackslash}p{1.8cm} >{\raggedright\arraybackslash}p{2.2cm} >{\raggedright\arraybackslash}p{2.8cm} >{\raggedright\arraybackslash}p{2.2cm}}
\toprule
\textbf{Run ID} & \textbf{Parallel.} & \textbf{Attention} & \textbf{Optim.} & \textbf{Held Constant} & \textbf{Expected Outcome} & \textbf{Metrics} \\
\midrule
\endfirsthead
\toprule
\textbf{Run ID} & \textbf{Parallel.} & \textbf{Attention} & \textbf{Optim.} & \textbf{Held Constant} & \textbf{Expected Outcome} & \textbf{Metrics} \\
\midrule
\endhead
\bottomrule
\endfoot

R1 & Single GPU & Standard & None & ViT-Small, CIFAR-10, batch 64, 100 steps, seed 42, fp32 & Absolute reference: slowest configuration & Time/step, throughput \\
\midrule
R2 & Single GPU & Flash & None & Same as R1 & Faster than R1; maximum FA speedup & Time/step, throughput, speedup \\
\midrule
R3 & DDP (2 GPU) & Standard & None & Same as R1 & $\sim$1.5--1.8$\times$ speedup over R1 & Time/step, throughput, speedup, efficiency \\
\midrule
R4 & DDP (2 GPU) & Flash & None & Same as R1 & Faster than R3, FA speedup \% smaller than R2 vs R1 & Time/step, throughput, speedup, FA rel. speedup \\
\midrule
R5 & FSDP (2 GPU) & Standard & None & Same as R1 & Slower than R3 due to higher comm & Time/step, throughput, speedup, efficiency \\
\midrule
R6 & FSDP (2 GPU) & Flash & None & Same as R1 & Faster than R5, FA speedup \% smallest & Time/step, throughput, speedup, FA rel. speedup \\
\midrule
R7 & DDP (2 GPU) & Flash & AMP & Same as R1 (except precision) & Faster than R4; AMP stacks with FA & Time/step, throughput, AMP impact \\
\midrule
R8 & DDP (2 GPU) & Flash & Overlap & Same as R1 & Faster than R4; overlap benefit may be smaller with FA & Time/step, throughput, comm breakdown \\
\midrule
R9 & DDP (2 GPU) & Flash & AMP + Overlap & Same as R1 (except precision) & Fastest DDP config & Time/step, throughput, combined impact \\

\end{longtable}
\end{singlespace}

\vspace{1em}
\textbf{Key Comparisons:}
\begin{itemize}
    \item \textbf{Hypothesis test:} Compare FA relative speedup across (R2 vs R1) $\rightarrow$ (R4 vs R3) $\rightarrow$ (R6 vs R5). Expect: decreasing \%.
    \item \textbf{Optimization stacking:} Compare R4 $\rightarrow$ R7 (AMP effect), R4 $\rightarrow$ R8 (overlap effect), R4 $\rightarrow$ R9 (combined).
    \item \textbf{Comm overhead visualization:} Compute vs comm breakdown for R3, R4, R5, R6.
\end{itemize}

\vspace{1em}
\textbf{Plots to Produce:}
\begin{enumerate}
    \item Bar chart: time/step for R1--R6 (core hypothesis test)
    \item Grouped bar chart: FlashAttention relative speedup \% across Single GPU / DDP / FSDP
    \item Stacked bar chart: compute vs communication breakdown for R3--R6
    \item Bar chart: optimization ablation impact (R4, R7, R8, R9)
\end{enumerate}

% ============================================================
% Section C: Config Pack
% ============================================================
\section{Config Pack}

Configuration files are stored in \texttt{configs/} as YAML files, one per Run ID. Each config specifies all parameters needed to reproduce the run. See \texttt{configs/r1\_single\_gpu\_standard.yaml} through \texttt{configs/r9\_ddp\_flash\_amp\_overlap.yaml}.

\vspace{1em}
Base parameters shared across all runs:

\begin{verbatim}
model:
  name: vit_small
  patch_size: 16
  embed_dim: 384
  num_heads: 6
  num_layers: 12

dataset:
  name: cifar10
  num_classes: 10

training:
  batch_size_per_gpu: 64
  num_steps: 100
  warmup_steps: 10
  seed: 42
  optimizer: adamw
  learning_rate: 0.001
\end{verbatim}

Per-run overrides are documented in individual config files. See \texttt{README.md} for exact reproduction commands.

\end{document}
